% Quiz 07 --- answer key with marking notes. DO NOT HAND OUT.
%
% Every number was checked with numpy: the five-person chain (opinions after one,
% two, three rounds; the limit by P^200; the eigenvalues) and the knight on
% 8 x 8 (64 x 64 matrix, 336 = sum of degrees, eigenvalue range, the +-1 vector,
% the probability of the start colour at steps 1..4, the lazy matrix).
%
% Compile with: xelatex -interaction=nonstopmode solutions.tex   (run twice)
\documentclass[a4paper, 11pt]{extarticle}
% Shared preamble for Quiz 07 (quiz07.tex and solutions.tex).
% Compile with: xelatex -interaction=nonstopmode <file>.tex   (run twice)
%
% Same house look as the pen-and-paper sheets in
% lecture-note/*/pen-and-paper/: Charter, A4, generous leading, room to draw.

\usepackage[top=0.7in, bottom=0.6in, left=0.8in, right=0.8in]{geometry}
\usepackage{amsmath}
\usepackage{amssymb}
\usepackage{array}
\usepackage{graphicx}
\usepackage[hidelinks]{hyperref}
\usepackage{tcolorbox}
\usepackage{enumitem}
\usepackage{etoolbox}
\usepackage{tikz}
\usepackage{fontspec}
\usetikzlibrary{calc,positioning,arrows.meta}
\definecolor{pltblue}{HTML}{1F77B4}
\definecolor{rulegrey}{HTML}{6A6D75}

% Body face. Charter, and the four ways to get it. XCharter *is* Charter and
% ships with a full TeX Live, so it is asked for first and by file: a CI runner
% answers yes to \IfFontExistsTF{Charter}, then cannot typeset it.
\IfFontExistsTF{XCharter-Roman.otf}{
  \setmainfont{XCharter}[
    Extension      = .otf,
    UprightFont    = *-Roman,
    BoldFont       = *-Bold,
    ItalicFont     = *-Italic,
    BoldItalicFont = *-BoldItalic]
}{
  \IfFontExistsTF{Charter}{\setmainfont{Charter}}{
    \IfFontExistsTF{Iowan Old Style}{\setmainfont{Iowan Old Style}}{
      \IfFontExistsTF{Georgia}{\setmainfont{Georgia}}{}}}}

\setlength{\parindent}{0pt}
\setlength{\parskip}{0.4em}
\setlength{\emergencystretch}{1.2em}
\raggedbottom
\pagestyle{empty}

% ---------------------------------------------------------------------------
% The submission link.
%
% Both the QR square and the printed address are the course's own short link,
% not the form's: paper outlives any one form. Caddy on the course droplet
% points go.skojaku.com/ans-quiz07 at the current Google Form
% (/etc/caddy/Caddyfile, block go.skojaku.com). Swapping the form is a line on
% the server, not a reprint.
%
% If the link ever changes, rebuild the square with
%
%   uvx --from segno segno --output=quiz07-form-qr.png --scale=20 --border=1 \
%       --error=m "<the url>"
% ---------------------------------------------------------------------------
\newcommand{\formlink}{https://go.skojaku.com/ans-quiz07}
\newcommand{\formshort}{\href{\formlink}{\texttt{go.skojaku.com/ans-quiz07}}}

% Quiz header. #1 = title, #2 = the line under it, #3 = term.
\newcommand{\quizhead}[3]{%
  {\footnotesize\textsc{Advanced Network Science} \hfill \textsc{#3}}\par
  \vspace{-0.2em}\rule{\textwidth}{1pt}\par
  \vspace{0.3em}
  {\LARGE\bfseries #1}\par
  \vspace{0.1em}
  {\small #2}\par
  \vspace{0.7em}
  {\small Name \underline{\hspace{6cm}} \hfill Date \underline{\hspace{3.2cm}}}\par
  \vspace{0.7em}}

% A question heading with its point value.
\newcommand{\question}[3]{%
  \vspace{0.5em}
  {\large\bfseries Question #1 \quad\normalsize\mdseries (#2 points)}\par
  \vspace{-0.3em}{\color{rulegrey}\rule{\textwidth}{0.6pt}}\par
  \vspace{0.1em}
  {\bfseries #3}\par
  \vspace{0.2em}}

% An empty box to draw or write in. #1 = height, #2 = caption under the rule.
\newcommand{\drawbox}[2]{%
  \begin{tcolorbox}[colback=white, colframe=rulegrey, boxrule=0.5pt,
                    sharp corners, height=#1, valign=top,
                    boxsep=2pt, left=4pt, right=4pt, top=3pt, bottom=3pt]
  {\footnotesize\color{rulegrey} #2}
  \end{tcolorbox}}

% Ruled writing lines, #1 of them.
\newcommand{\writelines}[1]{%
  \vspace{0.3em}
  \foreach \i in {1,...,#1}{%
    {\color{rulegrey}\rule{\textwidth}{0.4pt}}\par\vspace{1.5em}}}

\newgeometry{top=0.4in, bottom=0.4in, left=0.75in, right=0.75in}

\renewcommand{\question}[3]{\vspace{0.3em}
  {\large\bfseries Question #1 \quad\normalsize\mdseries (#2 points) --- #3}\par
  \vspace{-0.3em}{\color{rulegrey}\rule{\textwidth}{0.6pt}}\par\vspace{0.1em}}
\newcommand{\ans}[1]{{\color{pltblue}#1}}

\begin{document}

\quizhead{Quiz 7 \; --- \; Answer key}
        {Instructor copy \quad$\cdot$\quad do not hand out}
        {Fall 2026}

\question{1}{6}{Opinions on a network.}

Degrees: $k=(2,2,3,2,1)$, so $2m=10$. $P$ has rows
A: $(0,\tfrac12,\tfrac12,0,0)$, B: $(\tfrac12,0,\tfrac12,0,0)$,
C: $(\tfrac13,\tfrac13,0,\tfrac13,0)$, D: $(0,0,\tfrac12,0,\tfrac12)$,
E: $(0,0,0,1,0)$.

\textbf{(a) [1]} One round:
A $=(0.2+0.6)/2=\ans{0.4}$, B $=(0.8+0.6)/2=\ans{0.7}$,
C $=(0.8+0.2+0.4)/3=\ans{0.467}$, D $=(0.6+1.0)/2=\ans{0.8}$, E $=\ans{0.4}$.
The sum is $2.767$, not the starting $3.0$: \ans{no}. (1 point for the five
numbers and the answer to the question; $\tfrac12$ if the numbers are right and
the question is not answered.)

\textbf{(b) [2]} $w_i=\ans{k_i}$, or any multiple: $\ans{(0.2,\,0.2,\,0.3,\,0.2,\,0.1)}$.
Proof (1 point for the numbers, 1 for the proof):
\[
  \sum_i k_i\,x_i(t+1)=\sum_i k_i\sum_j \frac{A_{ij}}{k_i}\,x_j(t)
  =\sum_j\Bigl(\sum_i A_{ij}\Bigr)x_j(t)=\sum_j k_j\,x_j(t).
\]
In vector form: $\pi^\top P=\pi^\top$ for $\pi_i=k_i/2m$, so $\pi^\top x(t+1)=\pi^\top P\,x(t)=\pi^\top x(t)$.
\textbf{This is the left eigenvector} of $P$ for eigenvalue 1; $P$ acts on opinions
from the right, and the conserved quantity is the one that $P$ acts on from the
left. Check: after one round $0.2\cdot0.4+0.2\cdot0.7+0.3\cdot0.467+0.2\cdot0.8+0.1\cdot0.4=0.56$.
Full credit for an answer that finds the weights by guessing a pattern and
verifies the identity above; $\tfrac12$ for weights found and checked on the one
round of (a) only.

\textbf{(c) [1]} Everyone ends at the same $x^*$, so $\sum_i w_i x_i=x^*$, and
$x^*=\pi^\top x(0)=0.2\cdot0.8+0.2\cdot0.2+0.3\cdot0.6+0.2\cdot0.4+0.1\cdot1.0=\ans{0.56}$.
The plain mean of the starting opinions, $0.6$, is the common wrong answer: it
treats E, who has one friend and whom nobody but D listens to, like C, who has
three. $0.6$ gets 0.

\textbf{(d) [2]} \ans{$-0.857$} (1 point), so the factor is about \ans{0.857} per round
(the error also changes sign each round). It is the eigenvalue with the largest
\emph{absolute value} below 1, not the second largest eigenvalue $0.654$ (1 point
for the reason: $x(t)-x^*\mathbf 1$ is a sum of terms $\lambda_j^t\,v_j$ over the
other eigenvalues, and the term with the largest $|\lambda_j|$ decays last).
$0.654$: 0 for the pair. Check: the largest error after 5, 10, 20 rounds is
$0.108$, $0.055$, $0.011$.

% ---------------------------------------------------------------------------
\question{2}{4}{A knight on a chessboard.}

\textbf{(a) [1]} A jump changes column $+$ row by $\pm1\pm2$ or $\pm2\pm1$, an odd
number, so every jump changes the colour of the square. Start on a1 (dark):
dark with probability $\ans{0,\,1,\,0,\,1}$ for $t=1,2,3,4$. This does not
approach $\tfrac12$, so $p(t)$ \ans{does not converge} to $\pi$. (Both parts for the point.)

\textbf{(b) [2]} The vector $s$ with $s_i=+1$ on dark squares and $-1$ on light
squares (1 point). Every square the knight can reach from $i$ has the other
colour, so
\[
  (Ps)_i=\frac1{k_i}\sum_{j\sim i}s_j=\frac1{k_i}\,k_i\,(-s_i)=-s_i
\]
(1 point). So $-1$ is an eigenvalue. (A student who proves it as ``the graph is
bipartite, hence $-1$'' without the vector gets 1.)

\textbf{(c) [1]} The eigenvalues of $\tfrac12(I+P)$ are $\tfrac12(1+\lambda)$, so
$-1$ becomes \ans{$0$}, and every eigenvalue lies in $[0,1]$. Yes, $p(t)$ now
converges to $\pi$ (the stationary distribution of $\tfrac12(I+P)$ is the same
$\pi$). The convergence factor is then the second largest eigenvalue, which is
$\tfrac12(1+0.720)=0.860$ for the knight (not asked).

\vspace{0.5em}
\textbf{Facts, for the lecturer.} The knight's graph has 168 edges, so
$2m=336$; the dark and light squares each have total degree 168, which is why
$\pi$ gives $\tfrac12$ to each colour. A corner has 2 moves, a central square 8.
Degrees: 2 (4 squares), 3 (8), 4 (20), 6 (16), 8 (16).

\end{document}
